\chapter{A Dense Factorization of the Effective Reduced Hessian}
\chaptermark{Dense Factorization}
\label{chap:dense}

\vspace*{\fill}

\Cref{chap:updates} expressed every change of the working set as a row addition, a row
deletion or a rank-one modification of an $LDL^\top$ factorization, and realized these
operations with a sparse factorization of fixed dimension $\nmax$. This is the natural choice
for a matrix that is embedded in an index range of the size of the column capacity. The matrix
that actually has to be factorized, however, is small and dense: $\Heff$ never has more rows
than there are linking rows, and almost all of its entries are nonzero.

\medskip

This chapter develops a \textbf{dense factorization} that exploits these two properties. It
stores only the active indices, in the order in which they entered $\Nx$, so that a new index
is always \textbf{appended} as the last row of the factor. A removed index is deleted from the
factor by a \textbf{rank-one update} of the rows behind it, and the factor is compacted at the
same time. We derive these operations from the classical rank-one modification algorithm and
show four properties that make the dense factor attractive for the active-set method: the new
pivot of a released column is exactly its Schur complement, so the singularity test of
\cref{chap:singular} reduces to a single number; removals and updates never decrease a pivot;
downdates fail exactly when the modified matrix is not positive definite; and the factor never
contains a stored entry that does not belong to the current matrix, so that the periodic
refactorization of \cref{sec:refactorization} becomes unnecessary.

\medskip

The two factorizations are interchangeable: the active-set method only requires the
operations of \cref{sec:ldl-updates}, and both provide them. The chapter closes with a
comparison of their costs; the measurements follow in \cref{sec:results-factorizations}.

\vspace*{\fill}

\newpage


\section{Why a Dense Factor}
\label{sec:why-dense}

Three observations from the previous chapters motivate a different representation of the
factor.

\paragraph{The factorized matrix is small.} By \cref{cor:dimension}, as long as the reduced
Hessian is positive definite,
\begin{equation} \label{eq:dense-size}
|\Nx|\;\le\;\rank(\Aeff)\;\le\;m .
\end{equation}
The release of a dependent column (\cref{prop:singular-cases}) adds one more index before it is
detected and undone, so at most $m+1$ indices are ever active at the same time. The dimension of
$\Heff$ is therefore bounded by the number of linking rows and does not depend on the column
capacity $\nmax$ or on the number of generated columns. The fixed-length representation of
\cref{sec:fixed-length}, in contrast, gives the factor the dimension $\nmax$, which is typically
one to two orders of magnitude larger than $m$.

\paragraph{The factorized matrix is dense.} The entry $(\Heff)_{jj'} = u_j^\top\Meff u_{j'}$
of~\eqref{eq:implicit} is nonzero as soon as the differences $u_j = a_j - a_{\rep(j)}$ and
$u_{j'}$ share a row of $\Aeff$. Two columns of the same block share their representative, so
\[
(\Heff)_{jj'} = a_j^\top\Meff a_{j'} - a_j^\top\Meff a_{S_k} - a_{S_k}^\top\Meff a_{j'}
+ a_{S_k}^\top\Meff a_{S_k}
\qquad(\psi(j) = \psi(j') = k),
\]
which vanishes only by cancellation. Columns of different blocks compete for the same linking
rows, which is the purpose of these rows, and therefore overlap as well. The measurements of
\cref{sec:results-fill} confirm that the factor of $\Heff$ is practically dense, in agreement
with \cref{rem:natural-ordering}. A sparse factorization of a dense matrix stores the same
numbers as a dense one, together with row indices and column pointers, and its operations
follow these index structures instead of streaming through contiguous memory.

\paragraph{Sparse updates leave structural zeros.} As shown in \cref{sec:refactorization}, row
deletions and rank-one modifications of a sparse factor keep the positions of the removed
entries in its sparsity pattern. Every later operation pays for them until the factor is
rebuilt, and the refactorization interval has to balance the two costs.

\medskip

A dense factor of the active block only has none of these drawbacks, provided that it can be
updated as cheaply as the sparse one. The rest of this chapter shows that it can.


\section{Compact Representation}
\label{sec:compact}

Throughout this chapter, $t = |\Nx|$ is the number of active indices, and
$\Nx = \{j_1,\dots,j_t\}$ lists them in the order in which they entered $\Nx$: $j_1$ has been
active the longest and $j_t$ is the most recent. The \textbf{position} of an active index is
$\operatorname{pos}(j_\ell) = \ell$.

\begin{definition}[Compact representation] \label{def:compact}
The \textbf{compact effective reduced Hessian} is the $t\times t$ matrix
\[
\hat H_{\ell\ell'} = (\Heff)_{j_\ell j_{\ell'}}\qquad(\ell,\ell' = 1,\dots,t),
\]
i.e.\ $\Heff$ with its rows and columns arranged in activation order. The \textbf{dense
factor} is its $LDL^\top$ factorization without pivoting,
\[
\hat H = LDL^\top,\qquad L\in\R^{t\times t}\ \text{unit lower triangular},\qquad
D = \diag(d_1,\dots,d_t),
\]
stored as a dense array together with the list $j_1,\dots,j_t$ and the position map
$\operatorname{pos}$.
\end{definition}

The factor exists and is unique whenever $\Heff$ is positive definite, and then all pivots
$d_\ell$ are positive \cite{golub2013matrix}. The only object of size $\nmax$ is the integer
array $\operatorname{pos}$, with the value $0$ for inactive indices; the floating-point data
occupy at most $(m+1)^2$ entries by~\eqref{eq:dense-size}. The compact representation solves
the same systems as the embedded one:

\begin{lemma}[Solves in compact form] \label{lem:compact-solve}
Let $\zeta\in\R^{\nmax}$ vanish outside $\Nx$, and let $\hat\zeta_\ell = \zeta_{j_\ell}$. If
$\hat\eta$ solves $\hat H\hat\eta = \hat\zeta$, then $\eta$ with $\eta_{j_\ell} = \hat\eta_\ell$ and
$\eta_j = 0$ for $j\notin\Nx$ solves $\Heff^{\mathrm{emb}}\eta = \zeta$.
\end{lemma}

\begin{proof}
By \cref{lem:embedding}, the solution of $\Heff^{\mathrm{emb}}\eta = \zeta$ vanishes outside
$\Nx$, and its restriction to $\Nx$ solves $\Heff\eta_x = \zeta_{\Nx}$. The matrix $\hat H$ is $\Heff$
with rows and columns permuted by the same permutation as the entries of $\hat\zeta$ and
$\hat\eta$, so $\hat H\hat\eta = \hat\zeta$ is this system written in activation order.
\end{proof}

A solve therefore consists of a gather of $\zeta$ into activation order, a forward substitution
with $L$, a division by the pivots, a backward substitution with $L^\top$ and a scatter of the
result, at the cost of $2t^2 + O(t)$ arithmetic operations, independent of $\nmax$.

\medskip

The activation order is not a fill-reducing ordering, and it need not be: the matrix is dense,
so every ordering of its rows produces a dense factor. It is chosen because it turns a row
addition into the cheapest possible operation, an append.

\begin{center}
\renewcommand{\arraystretch}{1.25}
\begin{tabular}{lll}
\toprule
change of the working set & change of $\hat H$ & dense operation \\
\midrule
(i) column leaves $\Nx$ & position $\operatorname{pos}(q)$ removed & delete (\cref{sec:dense-delete}) \\
(ii) representative swap & rows of $\mathcal{N}_x^{(k)}$ rewritten & deletes + appends \\
(iii) surplus becomes active & $+\,\hat v_i\hat v_i^\top$ & rank-one update (\cref{sec:dense-rank-one}) \\
(iv) column enters $\Nx$ & new last row and column & append (\cref{sec:dense-append}) \\
(v) surplus becomes free & $-\,\hat v_i\hat v_i^\top$ & rank-one downdate (\cref{sec:dense-rank-one}) \\
\bottomrule
\end{tabular}
\end{center}

Here $\hat v_i$ is the vector $v_i$ of~\eqref{eq:rank-one-up} in activation order. A
representative swap is carried out as in \cref{sec:ldl-updates}: the replacement $S_k'$ is
deleted, and every row of $\mathcal{N}_x^{(k)}\setminus\{S_k'\}$ is deleted and appended again with
its new vector $z_j' = z_j - z_{S_k'}$. These rows thereby move to the end of the activation order.


\section{Appending an Index}
\label{sec:dense-append}

When a column $r$ is released (case~(iv) of \cref{thm:H-reduce}), $\Heff$ is bordered by the
vector $\gamma_r$ and the scalar $\gamma_{rr}$ of~\eqref{eq:bordered}. In the compact
representation the new index takes position $t+1$, so the new matrix is
\[
\hat H' = \begin{bmatrix}\hat H & \hat\gamma\\ \hat\gamma^\top & \gamma_{rr}\end{bmatrix},
\qquad \hat\gamma_\ell = (\gamma_r)_{j_\ell}.
\]

\begin{theorem}[Append] \label{thm:append}
Let $\hat H = LDL^\top$ be positive definite. Then
\[
\hat H' = \begin{bmatrix} L & 0\\ l^\top & 1\end{bmatrix}
\begin{bmatrix} D & 0\\ 0 & d_{t+1}\end{bmatrix}
\begin{bmatrix} L & 0\\ l^\top & 1\end{bmatrix}^\top
\qquad\text{with}\qquad
Lw = \hat\gamma,\quad l = D^{-1}w,\quad d_{t+1} = \gamma_{rr} - w^\top D^{-1}w .
\]
The factor of $\hat H$ is not changed, and the new row costs one forward substitution,
$t^2 + O(t)$ operations.
\end{theorem}

\begin{proof}
Multiplying out the right-hand side gives the blocks
\[
LDL^\top = \hat H,\qquad LDl = Lw = \hat\gamma,\qquad
l^\top Dl + d_{t+1} = w^\top D^{-1}w + d_{t+1} = \gamma_{rr},
\]
which are the blocks of $\hat H'$. The first factor is unit lower triangular and the middle one
diagonal, so this is the $LDL^\top$ factorization of $\hat H'$ whenever $d_{t+1}\ne0$, by
uniqueness. The vector $w$ is obtained from $Lw = \hat\gamma$ by forward substitution, and $l$
and $d_{t+1}$ in $O(t)$ further operations.
\end{proof}

The new pivot has a simple meaning. It is exactly the quantity that decides whether the
release is singular.

\begin{cor}[The new pivot is the Schur complement] \label{cor:append-pivot}
The new pivot of \cref{thm:append} is the Schur complement of~\eqref{eq:bordered},
\[
d_{t+1} = \sigma_r = \gamma_{rr} - \gamma_r^\top\Heff^{-1}\gamma_r
= \bigl\lVert (I - \Pi_{\mathrm{eff}})\,\Aeff z_r\bigr\rVert^2 ,
\]
where $\Pi_{\mathrm{eff}}$ is the orthogonal projector onto the range of $\Aeff Z_x$. In particular
$d_{t+1}\ge0$, and $d_{t+1} = 0$ if and only if the release of $r$ is singular.
\end{cor}

\begin{proof}
From $Lw = \hat\gamma$ and $\hat H = LDL^\top$,
\[
w^\top D^{-1}w = \hat\gamma^\top L^{-\top}D^{-1}L^{-1}\hat\gamma = \hat\gamma^\top\hat H^{-1}\hat\gamma
= \gamma_r^\top\Heff^{-1}\gamma_r ,
\]
since the permutation to activation order does not change this quadratic form. Hence
$d_{t+1} = \sigma_r$. With $U = \Aeff Z_x$ and $u = \Aeff z_r$ we have $\Heff = U^\top U$,
$\gamma_r = U^\top u$ and $\gamma_{rr} = u^\top u$, so
\[
\sigma_r = u^\top u - u^\top U(U^\top U)^{-1}U^\top u = u^\top(I - \Pi_{\mathrm{eff}})\,u
= \lVert(I-\Pi_{\mathrm{eff}})\,u\rVert^2,
\]
because $\Pi_{\mathrm{eff}} = U(U^\top U)^{-1}U^\top$ is symmetric and idempotent. The last statement is
part~(i) of \cref{prop:singular-cases}.
\end{proof}

\begin{remark}[Singularity detection with the dense factor] \label{rem:dense-detection}
\Cref{cor:append-pivot} makes the detection of \cref{sec:detection} exact for released columns.
An append leaves all existing pivots unchanged, and the new pivot, scaled by the pivot scale
$\lVert u_r\rVert^2$ of \cref{def:pivot-invariant}, is
\[
\frac{d_{t+1}}{\lVert u_r\rVert^2} = \frac{\lVert(I-\Pi_{\mathrm{eff}})\,\Aeff z_r\rVert^2}{\lVert u_r\rVert^2}
\;\le\;\sin^2\theta_r ,
\]
where $\theta_r$ is the angle between $\Aeff z_r$ and the range of $\Aeff Z_x$ (with equality if no
row of $u_r$ is masked). The pivot invariant thus tests directly whether the new direction is
numerically independent of the old ones. In the sparse factorization of \cref{prop:rowadd} the
same information is spread over the new pivot and the downdated pivots of the later indices.
A dependent column can moreover be removed again by deleting the last position, which by
\cref{thm:delete} below costs no arithmetic at all. (Our implementation undoes it by a
refactorization for both factorizations, which happens only a handful of times per run.)
\end{remark}


\section{Rank-One Modifications}
\label{sec:dense-rank-one}

The surplus cases~(iii) and~(v) change $\Heff$ by $\pm v_iv_i^\top$, and we will see in
\cref{sec:dense-delete} that the deletion of an index requires a rank-one update as well. All
three are instances of the following problem: given the factor $\hat H = LDL^\top$, a vector
$v\in\R^t$ and a scalar $\varphi$, compute the factor of $\hat H + \varphi\,vv^\top$. We use the
method C1 of Gill, Golub, Murray and Saunders \cite{gill1974methods}, and derive it here in a
form that also yields the properties we need.

\begin{theorem}[Rank-one modification] \label{thm:rank-one}
Let $\hat H = LDL^\top$ with $d_1,\dots,d_t>0$, let $v\in\R^t$ and $\varphi\in\R$, and suppose that
$\hat H + \varphi\,vv^\top$ is positive definite. Define $\varphi_1 = \varphi$, $v^{(0)} = v$, and
for $j = 1,\dots,t$
\begin{equation} \label{eq:c1}
\chi_j = v^{(j-1)}_j,\qquad \bar d_j = d_j + \varphi_j\chi_j^2,\qquad
\beta_j = \frac{\varphi_j\chi_j}{\bar d_j},\qquad \varphi_{j+1} = \frac{d_j}{\bar d_j}\,\varphi_j,\qquad
v^{(j)} = v^{(j-1)} - \chi_j\,Le_j .
\end{equation}
Then $\hat H + \varphi\,vv^\top = \bar L\bar D\bar L^\top$ with $\bar D = \diag(\bar d_1,\dots,\bar d_t)$ and the
unit lower triangular matrix $\bar L$ whose columns are
\[
\bar Le_j = Le_j + \beta_j\,v^{(j)}\qquad(j = 1,\dots,t).
\]
The vector $v^{(j)}$ vanishes in its first $j$ entries, so column $j$ of $\bar L$ differs from
column $j$ of $L$ only below the diagonal.
\end{theorem}

\begin{proof}
\emph{Step 1: reduction to a diagonal matrix.} Let $\tilde v = L^{-1}v$. Then
\[
\hat H + \varphi\,vv^\top = L\,(D + \varphi\,\tilde v\tilde v^\top)\,L^\top .
\]
If $D + \varphi\tilde v\tilde v^\top = \tilde L\bar D\tilde L^\top$, then
$\hat H + \varphi vv^\top = (L\tilde L)\,\bar D\,(L\tilde L)^\top$, and $\bar L = L\tilde L$ is unit lower triangular
as a product of two such matrices. It therefore suffices to factorize the diagonal matrix plus
a rank-one term.

\emph{Step 2: factor of $D + \varphi\tilde v\tilde v^\top$.} We show by induction over $t$ that
\begin{equation} \label{eq:c1-inner}
\bar d_j = d_j + \varphi_j\tilde v_j^2,\qquad \tilde L_{rj} = \beta_j\tilde v_r\;\;(r>j),\qquad
\beta_j = \varphi_j\tilde v_j/\bar d_j,\qquad \varphi_{j+1} = \varphi_jd_j/\bar d_j .
\end{equation}
Split off the first row and column:
\[
D + \varphi\tilde v\tilde v^\top =
\begin{bmatrix} d_1 + \varphi\tilde v_1^2 & \varphi\tilde v_1\tilde v_2^\top\\
\varphi\tilde v_1\tilde v_2 & D_2 + \varphi\tilde v_2\tilde v_2^\top\end{bmatrix},
\]
with $\tilde v = [\tilde v_1;\tilde v_2]$ and $D = \diag(d_1, D_2)$. One step of symmetric Gaussian
elimination gives the pivot $\bar d_1 = d_1 + \varphi\tilde v_1^2$, which is positive because the
matrix is positive definite, the first column $\varphi\tilde v_1\tilde v_2/\bar d_1 = \beta_1\tilde v_2$ of
$\tilde L$ below the diagonal, and the Schur complement
\[
D_2 + \varphi\tilde v_2\tilde v_2^\top - \frac{\varphi^2\tilde v_1^2}{\bar d_1}\tilde v_2\tilde v_2^\top
= D_2 + \varphi\Bigl(1 - \frac{\varphi\tilde v_1^2}{\bar d_1}\Bigr)\tilde v_2\tilde v_2^\top
= D_2 + \underbrace{\varphi\,\frac{d_1}{\bar d_1}}_{=\varphi_2}\,\tilde v_2\tilde v_2^\top .
\]
The Schur complement of a positive definite matrix is positive definite and has the same form
as the original matrix, with $D_2$, $\tilde v_2$ and $\varphi_2$, so the induction hypothesis
gives~\eqref{eq:c1-inner} for the remaining indices.

\emph{Step 3: the columns of $\bar L$.} By Step~2, column $j$ of $\tilde L$ is
$e_j + \beta_j\sum_{r>j}\tilde v_re_r$, so
\[
\bar Le_j = L\tilde Le_j = Le_j + \beta_j\sum_{r>j}\tilde v_r\,Le_r .
\]
Since $L\tilde v = v$, the sum equals $v - \sum_{i\le j}\tilde v_i\,Le_i$. We show that this is
$v^{(j)}$ and that $\chi_j = \tilde v_j$. Indeed, $v^{(j)} = v - \sum_{i\le j}\chi_i\,Le_i$ by its
recursion. If $\chi_i = \tilde v_i$ for $i<j$, then $v^{(j-1)} = \sum_{r\ge j}\tilde v_r\,Le_r$, whose
$j$-th entry is $\tilde v_j$ because $L$ is unit lower triangular; hence $\chi_j = \tilde v_j$. By
induction, $\chi_j = \tilde v_j$ for all $j$, and $\bar Le_j = Le_j + \beta_jv^{(j)}$. The same
representation shows that $v^{(j)} = \sum_{r>j}\tilde v_rLe_r$ vanishes in its first $j$ entries.
Inserting $\chi_j = \tilde v_j$ into~\eqref{eq:c1-inner} gives~\eqref{eq:c1}.
\end{proof}

The recursion~\eqref{eq:c1} computes $\tilde v = L^{-1}v$ by forward substitution and builds the
new factor in the same pass, column by column. \Cref{alg:c1} states it in the form in which it
is implemented: the vector $v^{(j)}$ overwrites $v$, and the new column overwrites the old one.

\begin{algobox}{Rank-One Modification $LDL^\top + \varphi\,vv^\top$ (Method C1)}{alg:c1}
\State \textbf{Input:} $L$, $d$ of $\hat H = LDL^\top$, vector $v$, scalar $\varphi$
\State $j_0\gets$ first index with $v_{j_0}\ne0$ \Comment{columns before $j_0$ are unchanged}
\For{$j = j_0,\dots,t$}
    \State $\chi\gets v_j$; \textbf{if} $\chi = 0$ \textbf{then continue} \Comment{column $j$ and $\varphi$ unchanged}
    \State $\bar d\gets d_j + \varphi\chi^2$
    \If{$\bar d\le0$} \textbf{return} \emph{not positive definite} \EndIf
    \State $\beta\gets\varphi\chi/\bar d$; \;$\varphi\gets\varphi\,d_j/\bar d$; \;$d_j\gets\bar d$
    \For{$r = j+1,\dots,t$}
        \State $v_r\gets v_r - \chi\,L_{rj}$; \;$L_{rj}\gets L_{rj} + \beta\,v_r$
    \EndFor
\EndFor
\end{algobox}

If $\chi_j = 0$, then $\bar d_j = d_j$, $\beta_j = 0$ and $\varphi_{j+1} = \varphi_j$, and the column is
unchanged; this justifies the skip in \cref{alg:c1}. The cost is $2(t - j_0)^2 + O(t)$
operations, where $j_0$ is the position of the first nonzero of $v$. The following properties
of~\eqref{eq:c1} are used below.

\begin{prop}[Pivots of a rank-one modification] \label{prop:c1-pivots}
In the situation of \cref{thm:rank-one}:
\begin{enumerate}
    \item If $\varphi>0$, then $\varphi_j>0$ and $\bar d_j\ge d_j$ for all $j$: an update never
          decreases a pivot.
    \item The pivots change by the factor
          \[
          \prod_{j=1}^{t}\frac{\bar d_j}{d_j} = 1 + \varphi\,v^\top\hat H^{-1}v .
          \]
    \item If $\varphi<0$ and the recursion~\eqref{eq:c1} is run without the assumption that
          $\hat H + \varphi vv^\top$ is positive definite, then all $\bar d_j$ are positive if and only if
          $\hat H + \varphi vv^\top$ is positive definite, i.e.\ if and only if
          $1 + \varphi\,v^\top\hat H^{-1}v>0$.
\end{enumerate}
\end{prop}

\begin{proof}
(i) By induction: if $\varphi_j>0$, then $\bar d_j = d_j + \varphi_j\chi_j^2\ge d_j>0$ and
$\varphi_{j+1} = \varphi_jd_j/\bar d_j>0$.

(ii) The determinant of an $LDL^\top$ factorization is the product of its pivots, so
\[
\prod_j\frac{\bar d_j}{d_j} = \frac{\det(\hat H + \varphi vv^\top)}{\det\hat H}
= 1 + \varphi\,v^\top\hat H^{-1}v
\]
by the matrix determinant lemma \cite{golub2013matrix}.

(iii) The recursion performs symmetric Gaussian elimination without pivoting on
$D + \varphi\tilde v\tilde v^\top$ (Step~2 of the proof of \cref{thm:rank-one}). Elimination without
pivoting produces positive pivots $\bar d_1,\dots,\bar d_t$ if and only if all leading principal
minors are positive, i.e.\ if and only if the matrix is positive definite, and
$D + \varphi\tilde v\tilde v^\top$ is positive definite if and only if
$\hat H + \varphi vv^\top = L(D+\varphi\tilde v\tilde v^\top)L^\top$ is. Since $\hat H$ is positive definite and
$\hat H + \varphi vv^\top$ differs from it by a rank-one term, the latter is positive definite if and
only if its determinant is positive (its smallest eigenvalue is the only one that can become
nonpositive, by interlacing), i.e.\ by~(ii) if and only if $1 + \varphi v^\top\hat H^{-1}v>0$.
\end{proof}

For the two surplus cases, $\varphi = \pm1$ and $v = \hat v_i$. An activated surplus ($\varphi = 1$)
can never fail and only increases pivots. A released surplus ($\varphi = -1$) produces positive
pivots exactly when $\omega_i = 1 - v_i^\top\Heff^{-1}v_i>0$, i.e.\ exactly when its release is not
singular (part~(ii) of \cref{prop:singular-cases}); the product of the pivot ratios is $\omega_i$.
\Cref{alg:c1} stops as soon as a pivot is not positive, and the method then proceeds as in
\cref{alg:singular}.

\begin{remark}[Numerical behaviour] \label{rem:c1-stability}
For $\varphi>0$ no cancellation occurs in $\bar d_j = d_j + \varphi_j\chi_j^2$, and method~C1 is
numerically stable \cite{gill1974methods}. For $\varphi<0$ the subtraction can cancel, and the
error grows when $\omega_i$ is small, i.e.\ when the downdated matrix is nearly singular. This is
precisely the situation that the pivot invariant of \cref{def:pivot-invariant} detects, and it
is then resolved by the zero-curvature step of \cref{chap:singular}. Away from it, the factor
stays accurate over long runs; \cref{sec:results-factorizations} measures its residual.
\end{remark}


\section{Removing an Index}
\label{sec:dense-delete}

When the index $q\in\Nx$ leaves the nullspace set (case~(i) of \cref{thm:H-expand}), its row
and column are removed from $\hat H$. Let $\ell = \operatorname{pos}(q)$ and partition the
positions into the groups $1 = \{1,\dots,\ell-1\}$, $2 = \{\ell\}$ and $3 = \{\ell+1,\dots,t\}$:
\[
L = \begin{bmatrix} L_{11} & & \\ l_{21}^\top & 1 & \\ L_{31} & l_{32} & L_{33}\end{bmatrix},\qquad
D = \diag(D_1, d_\ell, D_3).
\]

\begin{theorem}[Delete] \label{thm:delete}
Removing row and column $\ell$ from $\hat H = LDL^\top$ gives the matrix with the factor
\[
\begin{bmatrix} L_{11} & 0\\ L_{31} & \bar L_{33}\end{bmatrix},\qquad
\diag(D_1, \bar D_3),\qquad\text{where}\quad
\bar L_{33}\bar D_3\bar L_{33}^\top = L_{33}D_3L_{33}^\top + d_\ell\,l_{32}l_{32}^\top .
\]
The trailing factor is a rank-one update with $\varphi = d_\ell>0$, computed by \cref{alg:c1}. The
blocks $L_{11}$, $D_1$ and $L_{31}$ are kept.
\end{theorem}

\begin{proof}
Multiplying out $LDL^\top$ block by block gives, for the blocks that remain after the removal,
\[
\hat H_{11} = L_{11}D_1L_{11}^\top,\qquad \hat H_{31} = L_{31}D_1L_{11}^\top,\qquad
\hat H_{33} = L_{31}D_1L_{31}^\top + d_\ell\,l_{32}l_{32}^\top + L_{33}D_3L_{33}^\top .
\]
The proposed factor has the blocks $L_{11}D_1L_{11}^\top = \hat H_{11}$,
$L_{31}D_1L_{11}^\top = \hat H_{31}$ and
\[
L_{31}D_1L_{31}^\top + \bar L_{33}\bar D_3\bar L_{33}^\top
= L_{31}D_1L_{31}^\top + L_{33}D_3L_{33}^\top + d_\ell\,l_{32}l_{32}^\top = \hat H_{33}.
\]
It is unit lower triangular with a diagonal middle factor, so by uniqueness it is the $LDL^\top$
factorization of the reduced matrix. The trailing matrix $L_{33}D_3L_{33}^\top + d_\ell l_{32}l_{32}^\top$
is positive definite because $d_\ell>0$, so \cref{thm:rank-one} applies with $\varphi = d_\ell$.
\end{proof}

\begin{figure}[ht]
\centering
\begin{tikzpicture}[x=0.36cm, y=0.36cm, font=\small]
    % before: lower triangle of L, t = 10, deleted position ell = 4
    \fill[cblue!25] (0,10) -- (0,0) -- (10,0) -- cycle;
    \fill[cred!35] (3,0) rectangle (4,7);                % column ell below the diagonal (l_32)
    \fill[cred!35] (0,6) -- (0,7) -- (3,7) -- (3,6) -- cycle; % row ell (l_21)
    \fill[cgreen!35] (4,0) -- (4,6) -- (10,0) -- cycle;   % L_33
    \draw[thick] (0,10) -- (0,0) -- (10,0) -- cycle;
    \draw (3,0) -- (3,7); \draw (4,0) -- (4,6); \draw (0,6) -- (4,6); \draw (0,7) -- (3,7);
    \node at (1.2,3) {$L_{31}$};
    \node at (6.2,1.6) {$L_{33}$};
    \node[cred!70!black, above] at (3.5,7) {$\ell$};
    \node[below] at (5,-0.3) {$L$ before};
    % arrow
    \draw[-{Stealth[length=3mm]}, thick] (11.2,4) -- (14.8,4)
        node[midway, above, align=center] {\scriptsize delete $\ell$};
    % after: t = 9
    \begin{scope}[xshift=5.8cm]
    \fill[cblue!25] (0,9) -- (0,0) -- (9,0) -- cycle;
    \fill[cgreen!35] (3,0) -- (3,6) -- (9,0) -- cycle;
    \draw[thick] (0,9) -- (0,0) -- (9,0) -- cycle;
    \draw (3,0) -- (3,6); \draw (0,6) -- (3,6);
    \node at (1.2,3) {$L_{31}$};
    \node at (5.4,1.6) {$\bar L_{33}$};
    \node[below] at (4.5,-0.3) {$L$ after};
    \end{scope}
    % legend
    \begin{scope}[xshift=10.2cm, yshift=0.6cm]
    \fill[cred!35] (0,7) rectangle (0.7,7.7); \node[right] at (0.9,7.35) {\scriptsize removed: $l_{21}$, $d_\ell$, $l_{32}$};
    \fill[cgreen!35] (0,5.6) rectangle (0.7,6.3); \node[right] at (0.9,5.95) {\scriptsize updated by $d_\ell\,l_{32}l_{32}^\top$};
    \fill[cblue!25] (0,4.2) rectangle (0.7,4.9); \node[right] at (0.9,4.55) {\scriptsize kept, rows below $\ell$ shifted up};
    \end{scope}
\end{tikzpicture}
\caption{Deleting position $\ell$ from the dense factor (\cref{thm:delete}). The leading
columns are kept, the trailing triangle receives a rank-one update with the removed column
$l_{32}$ and the positive weight $d_\ell$, and all positions behind $\ell$ move up by one, so
that the factor stays contiguous.}
\label{fig:dense-delete}
\end{figure}

In an implementation, the update of the trailing block and the compaction are one pass:
\cref{alg:c1} reads column $\ell+s$ of the old trailing block and writes the updated column to
position $\ell+s-1$, which has been read completely in the previous step
(\cref{fig:dense-delete}). The rows $\ell+1,\dots,t$ of the leading columns are moved up by one
position, and the activation list and the position map are updated.

\begin{prop}[Cost of a deletion] \label{prop:delete-cost}
Deleting position $\ell$ costs $2(t-\ell)^2 + O(t)$ arithmetic operations for the rank-one update
and $O\bigl(\ell\,(t-\ell)\bigr)$ memory moves for the compaction. Deleting the most recent index
($\ell = t$) costs no arithmetic; deleting the oldest one ($\ell = 1$) costs as much as a full
rank-one update.
\end{prop}

\begin{proof}
The trailing block has $t-\ell$ positions, and \cref{alg:c1} on it costs $2(t-\ell)^2 + O(t-\ell)$
operations. The block $L_{31}$ has $t-\ell$ rows and $\ell-1$ columns, all of which are moved.
\end{proof}

The cost of a deletion therefore depends on how long the removed index has been active. In the
active-set method most removed indices have entered $\Nx$ recently: a released column moves
away from its bound (\cref{lem:descent}), and the columns that block are typically those with
small weights, which have been released last. The measurements of
\cref{sec:results-factorizations} show that a deletion touches about a quarter of the factor on
average.

\begin{cor}[Expansions preserve the pivot invariant] \label{cor:expand-pivots}
An expansion of type~(i) or~(iii) of \cref{lem:expand-ordering} does not decrease any pivot of
an index that stays in $\Nx$. In particular, if the pivot invariant of
\cref{def:pivot-invariant} holds before such an expansion, it holds after it.
\end{cor}

\begin{proof}
In case~(i), the pivots $D_1$ are kept and $\bar D_3$ results from a rank-one update with
$\varphi = d_\ell>0$, so $\bar D_3\ge D_3$ entrywise by part~(i) of \cref{prop:c1-pivots}; each pivot
moves with its index to the position one up. In case~(iii), the update $+\hat v_i\hat v_i^\top$ has
$\varphi = 1>0$. The pivot scales $\lVert u_j\rVert^2$ of the remaining indices do not change, since
their representatives do not change.
\end{proof}

In exact arithmetic the same holds for the sparse factorization, whose row deletion is the same
rank-one update (\cref{sec:ldl-updates}). \Cref{cor:expand-pivots} shows that the pivot
invariant can only be violated by a reduction, which is where \cref{chap:singular} checks it.


\section{Structural Zeros and Refactorization}
\label{sec:dense-no-zeros}

The dense factor stores exactly the $t(t-1)/2$ entries below the diagonal of the $t\times t$
array, the $t$ pivots, and the activation list. This is all that it ever stores:

\begin{prop}[No structural zeros] \label{prop:no-zeros}
After every sequence of appends, deletions and rank-one modifications, the leading
$t\times t$ block of the stored array is the factor $L$ of the current matrix $\hat H$, and no
stored entry belongs to a matrix that is no longer current.
\end{prop}

\begin{proof}
By induction over the operations. \Cref{thm:append} writes the new last row and pivot and keeps
the leading block. \Cref{thm:rank-one} overwrites the columns from $j_0$ on and keeps the
others, without changing the dimension. \Cref{thm:delete} keeps the leading columns, overwrites
the trailing triangle with $\bar L_{33}$, and moves everything behind position $\ell$ up by one,
so that the result again occupies exactly the leading $(t-1)\times(t-1)$ block. Entries outside
the leading block are never read, and the next append overwrites the entries of its row before
using them.
\end{proof}

The cost of every operation therefore depends only on the current $t$, and the factor never has
to be rebuilt to restore its efficiency: there is nothing like the growth of the stored pattern
of \cref{sec:refactorization}. A refactorization is only needed for numerical reasons, namely
when a pivot violates the invariant of \cref{def:pivot-invariant} or a downdate fails
(\cref{alg:singular}). Such a refactorization appends the active indices one by one as in
\cref{def:refactorization}, at the cost of about $t^3/3$ operations
(\cref{prop:refactor-cost}); afterwards their activation order is the increasing index order. An
optional periodic refactorization could be used to bound the accumulation of rounding errors,
but \cref{sec:results-factorizations} shows that it is not needed in practice.

\begin{remark}[Reset] \label{rem:dense-reset}
Resetting the factor to the identity only clears the position map of the active indices, in
$O(t)$. The stored entries need not be zeroed, because every append writes its whole row before
it is read.
\end{remark}


\section{Comparison with the Sparse Factorization}
\label{sec:dense-comparison}

\Cref{tab:dense-vs-sparse} compares the operations of the two factorizations. For the sparse
factorization, the costs are those of a factor without structural zeros, i.e.\ directly after a
refactorization; between refactorizations they grow with the stored pattern.

\begin{table}[ht]
\centering
\small
\renewcommand{\arraystretch}{1.3}
\begin{tabular}{>{\raggedright}p{0.2\textwidth}>{\raggedright}p{0.36\textwidth}>{\raggedright\arraybackslash}p{0.36\textwidth}}
\toprule
operation & sparse factor of $\Heff^{\mathrm{emb}}$ (\cref{chap:updates}) & dense factor of $\hat H$ (this chapter) \\
\midrule
dimension & $\nmax$, identity outside $\Nx$ & $t = |\Nx|\le m$ \\
release of a column, (iv) & row addition at index $r$: solve, new column, downdate of the later indices & append: one forward substitution, $t^2$ \\
removal of a column, (i) & row deletion: rank-one update of the later indices & delete: rank-one update of the later positions, $2(t-\ell)^2$ \\
surplus, (iii) and (v) & rank-one update or downdate along the elimination tree & \cref{alg:c1} from the first nonzero, at most $2t^2$ \\
solve & sparse solve on $\reach_L(\Nx)$ & two triangular solves, $2t^2$ \\
new pivot of a release & $D_{22}'\ge\sigma_r$; later pivots change & $d_{t+1} = \sigma_r$; no other pivot changes \\
stored entries & grow with every update (\cref{sec:refactorization}) & exactly $t(t-1)/2$ (\cref{prop:no-zeros}) \\
periodic refactorization & required, every $\tau$ updates & not required \\
floating-point memory & proportional to the stored pattern & at most $(m+1)^2$ \\
\bottomrule
\end{tabular}
\caption{The two factorizations of the effective reduced Hessian. Costs are arithmetic
operations; $\ell$ is the position of the removed index.}
\label{tab:dense-vs-sparse}
\end{table}

For a dense $\Heff$, both factorizations perform the same arithmetic up to constant factors. The
dense one does it on contiguous memory with simple loops and triangular solves, without index
lookups, without an $\nmax$-dimensional embedding and without structural zeros. On the other
hand, the dense factor costs $(m+1)^2$ floating-point numbers and $O(t^2)$ operations per update
regardless of the sparsity of $\Heff$. A sparse factor remains preferable when $\Heff$ is
genuinely sparse, for instance when the blocks have disjoint sets of linking rows, and $m$ is
large. For the crew scheduling instances of \cref{chap:results}, where $\Heff$ is dense, the
dense factorization is the better choice, as \cref{sec:results-factorizations} shows.

\medskip

Nothing in \cref{chap:qpas,chap:nullspace,chap:singular} depends on the choice. The
active-set method, the hot start of \cref{thm:hot-start} and the singularity treatment of
\cref{alg:singular} only use the operations of the table, so the two factorizations
produce the same sequence of iterates in exact arithmetic, and \cref{alg:hot-qpas} applies
unchanged, with the refactorization in its last line triggered only by the pivot invariant.
